\begingroup
\ArithmeticHead{234}{L}{ANALYSIS. BRIEFWECHSEL.}
\SourceLeaf{241}
\pdfbookmark[1]{Pfaff an Gauss, 8. Dec. 1800}{x1-correspondence-2}
\NumberedItem{2}
\begin{center}\large\bfseries\textsc{Pfaff} an \textsc{Gauss}.\end{center}
\begin{center}\{Theuerster Freund!\end{center}
\par\indent Ich sage Ihnen meinen herzlichen Dank für Ihren freundschaftlichen Brief, dessen Inhalt mir sehr interessant gewesen ist\quad .\quad .\quad .\quad .

\par\indent Ihr Problem vom dem lezten Gliede der Reihe
\[
a,\ b,\ \left[\frac{a+b}{2},\ \surd\frac{b(a+b)}{2},\ \frac{a+b+\surd2b(a+b)}{4}\right],\ldots
\]
\noindent hat mir recht wohl gefallen. Als ich gestern Mittag Ihr Paket erhielt, nahm ich mir vor, die Auflösung bis nach Beendigung eines Geschäfts zu verschieben, das mir jezt zur Last fällt\quad .\quad .\quad .\quad .\quad .\quad Indessen gestern Abend, als ich zu Haus kam, konnte ich mich nicht enthalten, das Problem näher anzusehen\quad .\quad .\quad .\quad .\quad .,\quad und fand auch sehr bald die Auflösung,\quad .\quad .\quad .\quad .\quad Es ist nehmlich
\[
f(a,b)=\frac{\surd(b^2-a^2)}{\operatorname{Arc}\cos\dfrac ab}\mathpunct{\raisebox{.5ex}{,}}
\]
\noindent oder (für $b<a$),
\[
=\frac{\surd(a^2-b^2)}{\log\left(\dfrac{a+\surd(a^2-b^2)}b\right)}\mathpunct{\raisebox{.5ex}{.}}
\]
\noindent Ist $a=1$, $b=\sec\omega$, so wird
\[
f(1,\sec\omega)=\frac{\operatorname{tang}\omega}{\omega}\mathpunct{\raisebox{.5ex}{,}}
\]
\noindent welches im Kreis durch eine einfache Construction dargestellt werden kann. Ferner ist
\[
\log\left(\frac{1+u}{1-u}\right)=\frac{2u}{f(1,\surd(1-u^2))}\mathpunct{\raisebox{.5ex}{.}}
\]
\noindent Für $a=1$, $b=\sin\omega$ ist
\[
f(1,\sin\omega)=\frac{\cos\omega}{\log\operatorname{tg}\tfrac12\omega}\mathpunct{\raisebox{.5ex}{.}}
\]
\noindent .\quad .\quad .\quad .\quad .\quad Für das versiegelte Billet und für das mir dadurch verschaffte
\clearpage
\ArithmeticHead{235}{R}{PFAFF AN GAUSS.}
\SourceLeaf{242}
\noindent Vergnügen bin ich Ihnen sehr verbunden. Ich schicke Ihnen das Billet hierbey zurück: ob ich gleich, wie natürlich, nicht zweifeln kann, dass unsere Auflösungen übereinstimmen, so würde es mich doch freuen, wenn Sie mir gelegentlich das Billet wieder zum Andenken zurückschickten.\quad .\quad .\quad .\quad .

\par\indent Fast hätte ich vergessen, Ihnen im voraus zu Ihrer Entdeckung Glück zu wünschen. Mir ist es besonders lieb, dass die h[öhere] Arithmetik mehr mit der Analysis in Verbindung gebracht ist, da ich durch frühe Verwöhnung immer mehr und lieber mit Buchstaben als mit Zahlen zu thun gehabt habe.

\par\indent Schenken Sie ferner Ihr freundschaftliches Andenken
\begin{flushright}Ihrem ergebensten Freund\end{flushright}
\noindent\begin{tabular*}{\linewidth}{@{\hspace{1.2em}}l@{\extracolsep{\fill}}r@{}}
Helmstedt&\textsc{J. F. Pfaff.}\\
8. Dec. 1800.&
\end{tabular*}
\begin{center}\textls[120]{Nachschrift.}\end{center}
\par\indent Zu meinem vorigen Brief an Sie, th[eurer] Fr[eund], folgt hiebey noch ein kleines Supplement, aus Veranlassung einer Bemerkung, die Ihnen gewiss auch nicht wird entgangen seyn.\quad .\quad .\quad .\quad .

\par\indent Wird bey der Reihe $a,b\ldots$ in Ihrem Gesez bloss das geändert, dass zuerst das geometrische, und dann das arithm. Mittel genommen wird, so ist
\[
f'(a,b)=\frac{\surd(ab-b^2)}{\operatorname{Arc}\cos\surd\dfrac ba}\mathpunct{\raisebox{.5ex}{,}}
\]
\noindent wie schon daraus folgt, dass $f'(a,b)=f(b,\surd ab)$. Wenn bey Ihrem \textls[120]{arithm.-}\linebreak\textls[120]{geometrischen Mittel} eine ähnliche Verwechslung gemacht wird, so ist $M'(a,b)=M(b,a)$: wie unmittelbar erhellt, weil dort auch das $3^{\mathrm{te}}$ und $4^{\mathrm{te}}$ Glied bloss verwechselt sind. Ich gestehe indessen, dass ich diese beyden leichten Bemerkungen nicht zuerst auf dem einfachsten Weg gemacht habe\textsuperscript{*)}.
\begin{center}C[etera] ut in literis.\end{center}
\begin{flushright}\textsc{J. Fr. Pf.}\}\end{flushright}
\par\vfill\noindent\rule{.25\linewidth}{.4pt}\par\smallskip
\begingroup\footnotesize\noindent\textsuperscript{*)}sonderbar, dass mir nicht gleich einfiel, dass bei dem arithm.-geom. M. die erwähnte Verwechslung gar keinen Unterschied macht, wo dann auch natürlich $M(b,a)=M(a,b)$.
\par\endgroup
\smallskip\noindent\small $30^{*}$
\clearpage
\endgroup
